% !TEX root = main.tex
\documentclass[12pt,a4paper]{ltjsarticle}

%% Fonts
\usepackage{lmodern}
\usepackage{luatexja-preset}

%% Layout
\usepackage[top=30mm, bottom=25mm, left=25mm, right=25mm]{geometry}
%% Section heading spacing — reduce vertical space before/after \section and \subsection.
%% Use the starred (*) versions to suppress paragraph indentation after the heading.
\usepackage{titlesec}
\titlespacing{\section}{0pt}{3.5ex plus 1ex minus 0.5ex}{2.3ex plus 0.5ex}
\titlespacing{\subsection}{0pt}{3.25ex plus 1ex minus 0.5ex}{1.5ex plus 0.5ex}

%% Color / graphics
\usepackage[svgnames,table]{xcolor}
\usepackage{graphicx}
\usepackage{float}

%% Diagrams / charts
\usepackage{tikz}
\usetikzlibrary{calc}
\usepackage{pgfplots}
\pgfplotsset{compat=1.18}

%% Source code listings
\usepackage{listings}
\renewcommand{\lstlistingname}{ソースコード}

%% Hyperlinks (load near last)
\usepackage{hyperref}

%% Header / footer
\usepackage{fancyhdr}

%% Captions
\usepackage[format=hang, font=small, labelfont=bf]{caption}
\usepackage{subcaption}

%% Mathematics
\usepackage{amsmath}
\usepackage{amssymb}

%% Professional table rules
\usepackage{booktabs}

%% URL
\usepackage{url}

%% List customization
\usepackage{enumitem}

%% Wide tables
\usepackage{tabularx}
\usepackage{multirow}

%% ── Source Code Appearance ───────────────────────────────────
\definecolor{codebg}{RGB}{248,248,248}
\definecolor{codeframe}{RGB}{180,180,180}
\definecolor{kwcolor}{RGB}{0,0,180}
\definecolor{cmcolor}{RGB}{100,100,100}
\definecolor{stcolor}{RGB}{160,0,0}

\lstset{
  backgroundcolor  = \color{codebg},
  basicstyle       = \ttfamily\small,
  keywordstyle     = \color{kwcolor}\bfseries,
  commentstyle     = \color{cmcolor}\itshape,
  stringstyle      = \color{stcolor},
  numbers          = left,
  numberstyle      = \tiny\color{gray},
  numbersep        = 8pt,
  frame            = single,
  rulecolor        = \color{codeframe},
  breaklines       = true,
  breakatwhitespace= false,
  showstringspaces = false,
  tabsize          = 4,
  captionpos       = b,
  xleftmargin      = 2em,
  framexleftmargin = 1.5em,
}

%% ── Page Style ───────────────────────────────────────────────
\pagestyle{fancy}
\fancyhf{}
\rhead{\small プログラミング実験}
\lhead{\small \nouppercase{\leftmark}}
\cfoot{\thepage}
\renewcommand{\headrulewidth}{0.4pt}

\fancypagestyle{titlepage}{
  \fancyhf{}
  \renewcommand{\headrulewidth}{0pt}
}
\fancypagestyle{frontmatter}{
  \fancyhf{}
  \cfoot{\thepage}
  \renewcommand{\headrulewidth}{0pt}
}

%% ── Hyperlink / PDF metadata ─────────────────────────────────
\hypersetup{
  colorlinks = true,
  linkcolor  = NavyBlue,
  citecolor  = NavyBlue,
  urlcolor   = NavyBlue,
  pdftitle   = {多次元データ処理},
  pdfauthor  = {Deka Risman Permana},
  pdfsubject = {プログラミング実験},
  pdfkeywords= {人間情報システム工学科 4年},
}

\newcommand{\HRule}{\rule{\linewidth}{0.5mm}}

\title{多次元データ処理}
\date{\today}
\author{Deka Risman Permana}

%% ============================================================
\begin{document}
%% ============================================================

%% ── Title Page ───────────────────────────────────────────────
\begin{titlepage}
  \thispagestyle{titlepage}
  \centering
  \vspace*{2cm}

  {\large 情報工学実験II}\\[0.5em]
  {\large 実験レポートレポート}\\[2em]

  \HRule{}\\[1em]
  {\LARGE \textbf{多次元データ処理}}\\[0.5em]
  \HRule{}

  \vspace{3cm}

  \begin{tabular}{rl}
    \textbf{氏　　名} & ：Deka Risman Permana \\[0.4em]
    \textbf{学　　科} & ：人間情報システム工学科 4年 \\[0.4em]
    \textbf{出席番号} & ：4347 \\[0.4em]
    \textbf{担当教員} & ：先生 \\
  \end{tabular}

  \vfill

  \textbf{提出日：}\today

\end{titlepage}

%% ── Table of Contents ────────────────────────────────────────
\newpage
\thispagestyle{frontmatter}
\pagenumbering{roman}
\tableofcontents

\newpage
\pagenumbering{arabic}
\setcounter{page}{1}

%% ── Sections ─────────────────────────────────────────────────

\section{実験目的}\label{sec:objective}
重要なデータ構造の 1 つである高階テンソル（多次元配列）の概要が理解でき、高階テンソルの
処理に関するプログラミング演習を通して、高階テンソルの生成、部分テンソル和の計算、行列展
開および畳み込みなどの操作ができるようになる。

\section{実験内容}\label{sec:content}
WebClass にアップされた教材の学習を行い、演習課題およびアンケートの回答などを実施する。


\subsection{演習I}\label{subsec:exercise1}

\subsubsection{実装}

\subsubsection{結果}

\subsubsection{説明}


\subsection{演習II}\label{subsec:exercise2}

\subsubsection{実装}

\subsubsection{結果}

\subsubsection{説明}


\subsection{演習III}\label{subsec:exercise3}

\subsubsection{実装}

\subsubsection{結果}

\subsubsection{説明}


\subsection{演習IV}\label{subsec:exercise4}

\subsubsection{実装}

\subsubsection{結果}

\subsubsection{説明}


\subsection{応用課題}\label{subsec:application}

\subsubsection{実装}

\subsubsection{結果}

\subsubsection{説明}


\section{感想}\label{sec:reflection}

\addcontentsline{toc}{section}{参考文献}
\begin{thebibliography}{9}
\end{thebibliography}

\end{document}
